% !TeX root = ../教材题库_全部条目_转写工作稿.tex
% Source: XJ-B1-C04-SU-P150-Y0796
% 第4章 / 小结与复习 小结与复习 / 教材 p.150 / PDF 第157页
% OCR未校对：公式与文字均可能有误；未完成数学TeX转写；禁止直接发给学生
% 原题及相邻原解截图（核对用）：../题目原图/XJ-B1-C04-SU-P150-Y0796-01.webp
% 本文件只在转写完成后作为题干使用；OCR 条目尚不是可靠题干。
\begin{bankitem}{第4章 小结与复习 第1题}{第4章 / 小结与复习 小结与复习 / 教材 p.150 / PDF 第157页}{XJ-B1-C04-SU-P150-Y0796}{OCR未校对：公式与文字均可能有误；未完成数学TeX转写；禁止直接发给学生}
\begin{ocrdraft}
[教材p.150 / PDF第157页]
\par
1.将整数指数幂推广到有理数指数幂和实数指数幂是描述客观世界所必需
\par
的，试描述指数幂的上述拓展过程，并说明我们是如何确定产的大小的？
\end{ocrdraft}
\end{bankitem}
